\documentclass[12pt]{article}
\usepackage{sbc-template}
\usepackage{graphicx,url}
\usepackage{booktabs}
\usepackage[utf8]{inputenc}
\usepackage{babel}
\usepackage{float}
\usepackage{multirow}
\usepackage{amsmath}
\usepackage{array}
\sloppy

\title{Desigualdade em Evidência: Uma Análise Socioeconômica Regional
do Desempenho no ENEM (2012--2024)}

\author{Tiago Defendi da Silva\inst{1}}

\address{Departamento de Computação -- DACOM\\
  Universidade Tecnológica Federal do Paraná -- UTFPR-CM\\
  Campo Mourão -- PR -- Brasil
  \email{tiagodefendi@alunos.utfpr.edu.br}
}

\begin{document}
\maketitle

\begin{resumo}
O Exame Nacional do Ensino Médio (ENEM) é o principal meio de acesso ao
ensino superior público no Brasil e a maior base longitudinal de dados
educacionais do país. Embora a literatura mostre que fatores
socioeconômicos afetam o desempenho discente, poucos estudos analisam se
a intensidade desse efeito varia entre macrorregiões ou incorporam
indicadores contextuais municipais, como o IDHM. Este trabalho investiga
em que medida o perfil socioeconômico dos candidatos e o IDHM explicam a
faixa de desempenho nas cinco macrorregiões do Brasil.
Foram usados os Microdados do ENEM de 2012 a 2024 (cerca de 51,3 milhões
de candidatos), integrados a indicadores do Atlas Brasil via DuckDB. Três
conjuntos de features foram comparados (M1: variáveis individuais; M2:
M1 + IDHM geral; M3: M1 + IDHM desagregado por cor/raça e gênero) usando
uma bateria progressivamente ampliada de classificadores: na abordagem
inicial, KNN, SVM e MLP com SMOTE; na abordagem estendida, Regressão
Logística, Árvore de Decisão, Random Forest, HistGradientBoosting, SVM
(LinearSVC), LightGBM e quatro configurações de ensemble, com
hiperparâmetros otimizados por Algoritmo Genético. Os experimentos foram
conduzidos sobre amostras de 1 M, 5 M e o dataset completo ($\approx$13 M
candidatos com IDHM disponível), usando divisão estratificada por ano
(2019--2024) e conjunto fixo de teste e validação. O melhor resultado foi
Macro F1 = 0,4893 (SVM LinearSVC, M3, 1 M amostras, limiares 500/700) e
Macro F1 = 0,4860 (Ensemble SVM+LGB, M2, dataset completo, limiares
500/730). A acurácia regional revelou que o Norte apresenta maior
facilidade de classificação (65,1\%), enquanto o Sudeste apresenta a menor
(57,0\%), resultado contraintuitivo que merece investigação adicional.
\end{resumo}

% ================================================================
\section{Introdução}
% ================================================================

O Exame Nacional do Ensino Médio (ENEM) é, desde sua reformulação em
2009, o principal mecanismo de acesso ao ensino superior público no
Brasil. Com mais de 75 milhões de inscrições acumuladas, o exame reúne
informações únicas sobre o desempenho de candidatos de todas as regiões
e perfis socioeconômicos.

A literatura indica que renda, tipo de escola, escolaridade dos pais,
raça e gênero estão sistematicamente associados às
notas~\cite{jaloto2021fatores,melo2021impacto}. Contudo, a maioria dos
estudos analisa apenas um ano, usa nota geral e não compara se o peso
desses fatores varia entre macrorregiões~\cite{dutra2023fatores}. Poucos
incorporam indicadores contextuais municipais como o IDHM em modelos
preditivos.

Este trabalho preenche essas lacunas integrando os microdados do ENEM
(2012--2024) com indicadores do Atlas Brasil (IDHM geral e desagregado
por cor/raça e gênero), treinando modelos de classificação por faixa de
nota em três versões e avaliando o ganho do IDHM. As perguntas de
pesquisa são:

\textbf{P1:} O perfil socioeconômico e o IDHM explicam a faixa de
desempenho no ENEM por área nas macrorregiões?

\textbf{P2:} O IDHM desagregado melhora os modelos em relação ao IDHM
geral ou às variáveis individuais?

% ================================================================
\section{Trabalhos Relacionados}
% ================================================================

Jaloto e Primi~\cite{jaloto2021fatores} analisam o ENEM 2018 com
regressão multinível, identificando nível socioeconômico como fator de
maior peso. Este trabalho amplia para 13 anos, adota classificação por
faixa e compara regiões.

Melo et al.~\cite{melo2021impacto} fazem análise espacial do ENEM 2018
integrada ao IBGE por município via LASSO e I de Moran. Este trabalho
opera em nível individual e usa IDHM desagregado como feature.

Ernica e Rodrigues~\cite{ernica2020desigualdades} investigam como NSE,
raça e gênero se combinam em metrópoles. Este trabalho aplica perspectiva
similar em escala nacional com IDHM desagregado como contexto externo.

Tertulino et al.~\cite{tertulino2026fatores} combinam Random Forest e
SHAP no ENEM 2022 no RN com seleção Boruta. Este trabalho estende para
escala nacional com ablation study M1/M2/M3.

Kleinke~\cite{kleinke2017influencia} analisa a influência do status
socioeconômico nos itens de Física do ENEM 2012, mostrando que o efeito
varia conforme o tipo de questão. Este trabalho amplia essa perspectiva
para Matemática, 13 anos de dados e comparação entre macrorregiões.

Dutra et al.~\cite{dutra2023fatores} revisam sistematicamente os fatores
que afetam o ENEM, apontando ausência de contexto municipal como lacuna.
Este trabalho responde diretamente a essa lacuna.

% ================================================================
\section{Metodologia}
% ================================================================

\subsection{Dados}

O dataset principal são os Microdados do ENEM 2012--2024~\cite{inep2024}
(76,3 milhões brutos; 51,3 milhões após filtro de presença completa nos
dois dias). A integração com o Atlas Brasil~\cite{pnud2017} foi feita via
DuckDB, utilizando código IBGE do município e ano como chaves. Na versão
atual, os dados do Atlas estão disponíveis apenas em nível estadual. A
variável de renda Q006 foi harmonizada entre os regimes pré e pós-2015
para cinco faixas comparáveis (A--E).

\subsection{Variável Alvo}

A variável alvo é a faixa de nota em Matemática, escolhida por apresentar
a maior disparidade regional (59 pontos entre Norte e Sudeste). Na
abordagem inicial utilizou-se a segmentação: $<$500 (Baixo), 500--700
(Médio) e $>$700 (Alto). Na abordagem estendida, os limiares foram
revisados para $<$500, 500--730 e $\geq$730, ajustando o ponto de corte
superior para refletir com maior precisão a distribuição empírica das
notas e a escassez real de candidatos com desempenho muito elevado.
Com os novos limiares, a distribuição no conjunto de treino é: Baixo
46,8\%, Médio 47,5\% e Alto 5,7\%, evidenciando o forte desbalanceamento
na classe Alto.

\subsection{Conjuntos de Features (Ablation Study)}

São comparados três conjuntos de features: \textbf{M1} com variáveis
individuais do candidato (renda harmonizada, escolaridade do pai e da
mãe, cor/raça, sexo, tipo de escola, macrorregião); \textbf{M2} com M1
acrescido do IDHM geral e seus subíndices (educação, renda, longevidade);
e \textbf{M3}, que substitui o IDHM geral de M2 pelo IDHM desagregado
por cor/raça (branco/negro) e por gênero (homem/mulher), mantendo as
demais variáveis de M1.

\subsection{Abordagem Inicial}

Para iniciar a modelagem, foram treinados quatro classificadores: KNN,
SVM (LinearSVC), LightGBM e MLP.

\textbf{KNN} classifica cada candidato pela classe majoritária entre os
$k$ vizinhos mais próximos. Os hiperparâmetros ($k$ e métrica de
distância) foram otimizados por busca em grade sobre $k \in
\{5,7,9,11,21\}$ e métrica $\in \{$euclidiana, manhattan$\}$.

\textbf{SVM (LinearSVC)} encontra a fronteira de decisão de margem
máxima. Os hiperparâmetros $C$ (em escala logarítmica, $\log_{10}(C) \in
[-3,2]$) e \textit{class\_weight} foram otimizados por Algoritmo
Genético (população de 25 indivíduos, 20 gerações, fitness calculado em
subconjunto de 30 mil amostras).

\textbf{MLP} (três camadas densas, 256→128→64 neurônios) com SMOTE
aplicado ao conjunto de treino para balancear as classes, rebalanceando
de 52,3\% / 40,6\% / 7,1\% para 33,3\% por classe.

\textbf{LightGBM} com hiperparâmetros otimizados por Algoritmo Genético
(440 estimadores, 115 folhas, taxa de aprendizado 0,0289, mínimo de 79
amostras por folha) e SMOTE sobre o mesmo subconjunto do MLP.

\textbf{Ensemble:} combinação das predições de KNN, SVM e LightGBM em
votação hard e soft, com a saída do SVM convertida em probabilidades via
softmax antes de combinar.

Os experimentos iniciais foram conduzidos sobre amostras de 500 mil a 1
milhão de registros, com divisão em 60\% treino, 20\% validação e 20\%
teste.

\subsection{Abordagem Estendida}

Diante dos resultados obtidos na abordagem inicial --- em especial o fraco
desempenho na classe Alto e a limitada capacidade dos modelos sem
técnicas de balanceamento ---, foi desenvolvida uma abordagem estendida
com as seguintes diferenças metodológicas:

\subsubsection{Engenharia de Features}

Todas as variáveis categóricas do questionário socioeconômico (colunas
\texttt{Q\_*}) foram sistematicamente recodificadas como valores
ordinais inteiros, armazenadas com sufixo \texttt{\_NUM}. O ano do exame
(\texttt{NU\_ANO}) foi incluído como feature numérica, permitindo que os
modelos capturem tendências interanuais no desempenho e nas condições
socioeconômicas. As features M1, M2 e M3 passaram a conter 28, 37 e 39
variáveis, respectivamente (Tabela~\ref{tab:features}).

\begin{table}[H]
\centering
\caption{Composição dos conjuntos de features na abordagem estendida}
\label{tab:features}
\small
\begin{tabular}{lcp{8cm}}
\toprule
\textbf{Conjunto} & \textbf{Qtd.} & \textbf{Conteúdo} \\
\midrule
M1 & 28 & Variáveis individuais: NU\_ANO, macrorregião, sexo, cor/raça,
           tipo de escola, escolaridade dos pais, renda, bens domésticos,
           situação de trabalho, bolsa família \\
M2 & 37 & M1 + IDHM geral, IDHM educação/renda/longevidade, renda per
           capita, tx.\ analfabetismo, tx.\ envelhecimento, esperança de
           vida, mortalidade infantil \\
M3 & 39 & M2 + IDHM desagregado por cor/raça do candidato
           (\texttt{idhm\_cand\_raca}) e por sexo
           (\texttt{idhm\_cand\_sexo}) \\
\bottomrule
\end{tabular}
\end{table}

O IDHM personalizado de M3 é derivado cruzando a raça do candidato com
o IDHM da raça correspondente no estado (\texttt{idhm\_branco} para
brancos, \texttt{idhm\_negro} para pretos e pardos) e o sexo com o IDHM
de gênero (\texttt{idhm\_homem}, \texttt{idhm\_mulher}), tornando o
indicador contextual sensível à identidade do próprio candidato.

\subsubsection{Divisão Estratificada por Ano}

Para garantir que todos os anos de 2019 a 2024 estejam proporcionalmente
representados em treino, validação e teste, adotou-se divisão
estratificada por ano: para cada ano, 10\% dos candidatos foram
reservados para o conjunto de teste e 10\% para validação, ambos
amostrados de forma estratificada pela classe alvo; os 80\% restantes
compõem o pool de treino. Os conjuntos de teste e validação foram salvos
em arquivos fixos (\texttt{test\_set.parquet} e \texttt{val\_set.parquet})
e reutilizados por todos os classificadores, garantindo comparação justa.
O conjunto de treino é amostrado do pool remanescente a cada experimento,
excluindo explicitamente as linhas presentes em teste e validação.

\subsubsection{Algoritmos}

A bateria de classificadores foi ampliada para cobrir um espectro mais
amplo de paradigmas de aprendizado:

\begin{itemize}
  \item \textbf{Regressão Logística (LR):} modelo linear com regularização
        L2 e solver SAGA, adequado a grandes volumes de dados.
  \item \textbf{Árvore de Decisão (DT):} baseline interpretável,
        profundidade máxima 12.
  \item \textbf{Random Forest (RF):} ensemble de 80 árvores com
        subamostragem de 10--30\% por árvore (\texttt{max\_samples}),
        necessária para viabilizar o treino no dataset completo.
  \item \textbf{HistGradientBoosting (HGB):} implementação sklearn de
        gradient boosting histogramado, nativa para dados com valores
        ausentes.
  \item \textbf{SVM (LinearSVC):} SVM linear otimizado para grandes
        volumes (\texttt{dual=auto}).
  \item \textbf{LightGBM (LGB):} gradient boosting eficiente, usado com
        300--440 estimadores.
\end{itemize}

Para os quatro modelos sensíveis a hiperparâmetros (LR, SVM, LGB, HGB),
os hiperparâmetros foram otimizados por Algoritmo Genético (AG): os genes
codificam o espaço de busca de cada modelo, a fitness é avaliada em um
subconjunto de 30--50 mil amostras de treino e 10--15 mil de validação, e
o melhor indivíduo é usado para treinar o modelo final no dataset
completo. O AG usa seleção por torneio ($k=3$), cruzamento uniforme
($p=0,70$), mutação gaussiana para genes contínuos e mutação de flip para
genes binários, com elitismo preservando os dois melhores indivíduos por
geração.

Foram construídos quatro ensembles:

\begin{itemize}
  \item \textbf{Ensemble SVM+LGB Hard:} votação majoritária entre SVM e
        LGB.
  \item \textbf{Ensemble SVM+LGB Soft:} média de probabilidades entre
        SVM (via SGDClassifier com \textit{modified\_huber loss}, que
        emite probabilidades nativamente) e LGB.
  \item \textbf{Ensemble Full Soft:} média de probabilidades de SGD-SVM,
        LGB, RF e HGB.
  \item \textbf{Ensemble Top-2 Soft:} ensemble dinâmico dos dois modelos
        base com maior Macro F1 no conjunto M1, treinado e avaliado
        também em M2 e M3.
\end{itemize}

\subsubsection{Escalabilidade: Múltiplos Tamanhos de Amostra}

Para investigar o efeito do volume de treino sobre o desempenho dos
classificadores, foram executados experimentos com tamanhos de amostra
crescentes: \textbf{1 M} de candidatos (excluindo teste e validação),
\textbf{5 M} e o \textbf{dataset completo} ($\approx$13,4 M candidatos
com IDHM disponível nos anos 2019--2024). A comparação foi conduzida
mantendo os mesmos conjuntos de teste e validação em todos os casos.

\subsection{Hipóteses}

\textbf{H1:} A correlação renda--desempenho é mais intensa no
Norte/Nordeste do que no Sul/Sudeste.
\textbf{H2:} O efeito socioeconômico é maior em Matemática e menor em
Redação.
\textbf{H3:} Modelos com IDHM (M2 e M3) superam em Macro F1 o modelo
baseline M1.

\subsection{Avaliação}

A métrica principal é o Macro F1, por ser robusta ao desbalanceamento de
classes e dar peso igual a todas as faixas, inclusive à minoritária
(Alto). A Acurácia é reportada como referência complementar. Todas as
métricas são calculadas sobre o conjunto de validação, exceto os
resultados finais (reservados para a entrega final).

% ================================================================
\section{Resultados e Discussão}
% ================================================================

\subsection{Disparidade Regional e Correlação Renda--Desempenho}

A Tabela~\ref{tab:medias} mostra as notas médias por macrorregião. A
disparidade Norte--Sudeste é de 59 pontos em Matemática e 32,9 em
Ciências da Natureza, padrão consistente em todos os anos da série,
suportando H1 visualmente.

\begin{table}[H]
\centering
\caption{Nota média por macrorregião e área (2012--2024)}
\label{tab:medias}
\small
\begin{tabular}{lccccc}
\toprule
\textbf{Região} & \textbf{CN} & \textbf{CH} & \textbf{LC} & \textbf{MT} & \textbf{Redação} \\
\midrule
Norte        & 464,6 & 504,8 & 485,4 & 473,8 & 520,5 \\
Nordeste     & 471,6 & 510,7 & 493,2 & 491,7 & 537,0 \\
Centro-Oeste & 485,4 & 527,5 & 508,5 & 506,6 & 548,4 \\
Sul          & 495,9 & 541,4 & 521,7 & 526,7 & 561,4 \\
Sudeste      & 497,5 & 544,0 & 525,5 & 532,8 & 570,9 \\
\midrule
\textbf{Brecha SE--N} & 32,9 & 39,2 & 40,1 & \textbf{59,0} & 50,4 \\
\bottomrule
\end{tabular}
\end{table}

As correlações de Spearman entre renda e nota são moderadas
($\rho \approx 0,45$ em Matemática) e variam por macrorregião, com o
Nordeste apresentando correlações maiores que o Sudeste, possivelmente
por maior heterogeneidade intra-regional.

\subsection{Resultados da Abordagem Inicial}

O KNN sobre M1 (Macro F1 = 0,4787) serve de baseline. A
Tabela~\ref{tab:resultados_inicial} resume os resultados iniciais.

\begin{table}[H]
\centering
\caption{Resultados da abordagem inicial -- Macro F1 e Accuracy (validação, FAIXA\_MT)}
\label{tab:resultados_inicial}
\small
\begin{tabular}{lccc}
\toprule
\textbf{Modelo} & \textbf{Features} & \textbf{Macro F1} & \textbf{Accuracy} \\
\midrule
\textbf{KNN (baseline), $k$=21, manhattan} & M1 & \textbf{0,4787} & 0,6141 \\
KNN, $k$=11, manhattan                     & M2 & 0,4697 & 0,6014 \\
SVM (LinearSVC) + AG                       & M1 & 0,4660 & 0,5403 \\
SVM (LinearSVC) + AG                       & M2 & 0,4717 & 0,5453 \\
SVM (LinearSVC) + AG, 5-CV                 & M2 & 0,4739 & 0,5463 \\
Ensemble Hard Voting (KNN+SVM+LGB)         & M3 & 0,5142 & 0,5866 \\
Ensemble Soft Voting (KNN+SVM+LGB)         & M3 & 0,5334 & 0,6031 \\
MLP + SMOTE                                & M3 & 0,5081 & 0,5816 \\
LightGBM + AG + SMOTE                      & M3 & 0,5406 & 0,6239 \\
\bottomrule
\end{tabular}
\end{table}

O melhor resultado inicial foi LightGBM + AG + SMOTE com Macro F1 =
0,5406, ganho de 0,0619 sobre o baseline KNN/M1. O SMOTE foi
determinante para elevar o recall da classe Alto de $\approx$0,00 para
$\approx$0,37--0,40, ao custo de queda no recall de Baixo (de 0,77 para
$\approx$0,60).

\subsection{Resultados da Abordagem Estendida}

\subsubsection{Comparação por Tamanho de Amostra (1 M, limiares 500/700)}

A Tabela~\ref{tab:pivot_1m} apresenta o ablation study M1/M2/M3 com
amostra de 1 M de candidatos, mostrando o Macro F1 no conjunto de
validação para todos os classificadores. O conjunto de teste foi selado
e não foi avaliado nesta etapa.

\begin{table}[H]
\centering
\caption{Macro F1 por modelo e variante de features (1 M amostras, limiares 500/700, validação)}
\label{tab:pivot_1m}
\small
\begin{tabular}{lccc}
\toprule
\textbf{Modelo} & \textbf{M1 (28)} & \textbf{M2 (37)} & \textbf{M3 (39)} \\
\midrule
\textbf{SVM (LinearSVC)}      & 0,4872 & 0,4892 & \textbf{0,4893} \\
HistGradientBoosting          & 0,4691 & 0,4717 & 0,4737 \\
LightGBM                      & 0,4682 & 0,4728 & 0,4733 \\
Ensemble SVM+LGB Hard         & 0,4626 & 0,4665 & 0,4671 \\
DecisionTree                  & 0,4523 & 0,4493 & 0,4436 \\
Ensemble Full Soft            & 0,4462 & 0,4488 & 0,4500 \\
Ensemble SVM+LGB Soft         & 0,4415 & 0,4448 & 0,4450 \\
Ensemble Top-2 Soft           & --     & 0,4448 & 0,4461 \\
LogisticRegression            & 0,4367 & 0,4397 & 0,4396 \\
RandomForest                  & 0,4313 & 0,4310 & 0,4312 \\
\bottomrule
\end{tabular}
\end{table}

Na abordagem estendida sem SMOTE, o SVM (LinearSVC) emerge como melhor
modelo em todas as variantes, atingindo Macro F1 = 0,4893 com M3. Há
um padrão consistente de melhora de M1 para M2 e de M2 para M3 na
maioria dos modelos, indicando que o IDHM agrega valor preditivo
incremental, ainda que modesto. O Random Forest apresentou os piores
resultados entre os modelos base, possivelmente por overfitting em
features categóricas recodificadas.

\subsubsection{Dataset Completo com Novos Limiares (500/730)}

Os experimentos com o dataset completo ($\approx$13,4 M candidatos de
treino) adotaram os limiares revisados de 500 e 730 pontos. A
Tabela~\ref{tab:resultados_full} apresenta o melhor modelo e o relatório
de classificação no conjunto de validação.

\begin{table}[H]
\centering
\caption{Melhor modelo -- dataset completo, limiares 500/730 (validação)}
\label{tab:resultados_full}
\small
\begin{tabular}{lcccc}
\toprule
\textbf{Classe} & \textbf{Precisão} & \textbf{Recall} & \textbf{F1} & \textbf{Suporte} \\
\midrule
Baixo ($<$500)       & 0,61 & 0,75 & 0,67 & 696.506 \\
Médio (500--730)     & 0,62 & 0,48 & 0,54 & 707.105 \\
Alto ($\geq$730)     & 0,23 & 0,26 & 0,24 &  85.160 \\
\midrule
\textbf{Macro avg}   & 0,49 & 0,49 & \textbf{0,49} & 1.488.771 \\
Weighted avg         & 0,59 & 0,59 & 0,59 & 1.488.771 \\
\textbf{Accuracy}    & \multicolumn{3}{c}{0,59} & 1.488.771 \\
\bottomrule
\end{tabular}
\end{table}

O melhor modelo no dataset completo foi o \textbf{Ensemble SVM+LGB Hard}
sobre M2 (Macro F1 = 0,4860), resultado próximo ao SVM/M3 com 1 M
amostras (0,4893), o que indica que o aumento de escala de 1 M para 13 M
não produziu ganho expressivo no Macro F1 com os modelos testados. O
desbalanceamento da classe Alto (5,7\% das amostras com os limiares
500/730) permanece como principal limitador, com F1 = 0,24 mesmo com 13 M
amostras de treino.

\subsubsection{Análise Regional}

A Tabela~\ref{tab:regional} apresenta a acurácia do melhor modelo
(Ensemble SVM+LGB Hard, M2) por macrorregião no conjunto de validação,
avaliando indiretamente H1.

\begin{table}[H]
\centering
\caption{Acurácia por macrorregião -- Ensemble SVM+LGB Hard (M2),
         dataset completo, limiares 500/730}
\label{tab:regional}
\small
\begin{tabular}{lcc}
\toprule
\textbf{Macrorregião} & \textbf{Acurácia} & \textbf{N (validação)} \\
\midrule
Norte        & \textbf{0,651} & 162.381 \\
Nordeste     & 0,602 & 537.742 \\
Centro-Oeste & 0,579 & 122.537 \\
Sul          & 0,576 & 161.739 \\
Sudeste      & 0,570 & 504.372 \\
\midrule
\textbf{Total} & 0,590 & 1.488.771 \\
\bottomrule
\end{tabular}
\end{table}

O resultado mais surpreendente é que o Norte apresenta a maior acurácia
de classificação (65,1\%), apesar de ter as menores notas médias
(Tabela~\ref{tab:medias}). Uma hipótese explicativa é que a distribuição
de notas no Norte é mais concentrada --- com menor variância intraclasse
---, tornando mais fácil para o modelo atribuir corretamente a faixa. Em
contraste, o Sudeste, apesar de ter as maiores médias, apresenta maior
heterogeneidade interna (candidatos de escolas públicas e privadas de
qualidade muito distinta no mesmo estado), o que dificulta a
classificação.

\subsection{Comparação das Abordagens}

A Tabela~\ref{tab:comparacao} consolida os melhores resultados de cada
abordagem.

\begin{table}[H]
\centering
\caption{Comparação entre abordagens (validação)}
\label{tab:comparacao}
\small
\begin{tabular}{lccccc}
\toprule
\textbf{Abordagem} & \textbf{Modelo} & \textbf{Features} & \textbf{N treino} & \textbf{Limiares} & \textbf{Macro F1} \\
\midrule
Inicial    & LGB + AG + SMOTE      & M3 & 277k  & 500/700 & 0,5406 \\
Estendida  & SVM (LinearSVC) + AG  & M3 & 1 M   & 500/700 & 0,4893 \\
Estendida  & Ensemble SVM+LGB Hard & M2 & 13 M  & 500/730 & 0,4860 \\
\bottomrule
\end{tabular}
\end{table}

O LightGBM com SMOTE da abordagem inicial ainda apresenta o melhor
Macro F1 global (0,5406), mas deve-se notar que este resultado foi obtido
num subconjunto de apenas 277 mil registros com SMOTE agressivo, que
pode inflar artificialmente as métricas por geração de amostras
sintéticas próximas às reais de treino. A abordagem estendida, sem SMOTE
e com divisão estratificada por ano e conjunto de teste fixo, oferece uma
avaliação mais rigorosa e reprodutível, e o Macro F1 $\approx$0,49 deve
ser interpretado como o resultado mais confiável disponível até o momento.

\subsection{Discussão das Hipóteses}

\textbf{H1} é parcialmente suportada pela análise exploratória: a brecha
de 59 pontos em Matemática entre Norte e Sudeste é consistente. Contudo,
a análise de classificação regional (Tabela~\ref{tab:regional}) revela
que fatores regionais adicionais --- não capturados pelas features atuais
--- influenciam a dificuldade de classificar os candidatos de diferentes
regiões.

\textbf{H3} é parcialmente confirmada: há ganho consistente de Macro F1
de M1 para M2 e de M2 para M3 na maioria dos modelos da abordagem
estendida (Tabela~\ref{tab:pivot_1m}), mas os ganhos são pequenos
($\Delta \leq 0,005$), sugerindo que o IDHM estadual tem poder
explicativo incremental limitado frente às variáveis individuais. Dados
municipais poderiam ampliar esse ganho.

\subsection{Limitações}

A principal limitação desta etapa é que o IDHM está disponível apenas em
nível estadual, reduzindo a granularidade das features de M2 e M3. A
classe Alto permanece fortemente desbalanceada (5,7\% com limiares
500/730), mesmo no dataset completo de 13 M registros. Os experimentos
com 5 M de amostras foram configurados mas seus resultados definitivos
aguardam execução completa para comparação direta com 1 M e dataset
completo. O conjunto de teste permanece selado para garantir a validade
da avaliação final.

% ================================================================
\section{Conclusão}
% ================================================================

Este trabalho apresentou um pipeline reprodutível para classificação da
faixa de desempenho no ENEM com 51,3 milhões de candidatos (2012--2024),
integrado ao Atlas Brasil via DuckDB. Duas abordagens foram desenvolvidas:
a inicial, com KNN, SVM, MLP e LightGBM com SMOTE, obteve Macro F1 =
0,5406 (LGB + AG + SMOTE, M3, 277 mil amostras); a estendida, com
engenharia de features aprimorada (NU\_ANO, codificação numérica),
divisão estratificada por ano, conjunto fixo de teste e validação, e uma
bateria de 10 modelos e ensembles com AG para otimização de
hiperparâmetros, alcançou Macro F1 = 0,4893 (SVM/M3, 1 M amostras) e
0,4860 (Ensemble SVM+LGB, M2, 13 M amostras), com avaliação mais
rigorosa e reprodutível.

A análise regional revelou que o Norte tem a maior acurácia de
classificação (65,1\%), resultado contraintuitivo frente às menores médias
da região, hipótese atribuída à menor dispersão intraclasse. O
desbalanceamento da classe Alto permanece o principal limitador dos
resultados.

Os próximos passos são: avaliar os modelos no conjunto de teste selado,
completar os experimentos com 5 M de amostras, investigar técnicas de
balanceamento compatíveis com datasets de grande escala (class weights,
threshold calibration) e obter dados do Atlas Brasil em nível municipal
para ampliar o poder explicativo do IDHM.

\bibliographystyle{plain}
\bibliography{referencias}

\end{document}
